A geometric space is defined by the relations among its points: distances, angles, and the behaviour of geodesics—paths that locally minimise distance. In Euclidean geometry, geodesics are governed by axioms such as the parallel postulate, which ensures that parallel lines never intersect and that triangle angles sum to 180 degrees. When these properties fail, the space is said to be curved.

Curvature quantifies how a space deviates from the rules of Euclidean geometry. Positive curvature causes initially parallel lines to converge, as on the surface of a sphere. Negative curvature causes them to diverge, as in a hyperbolic plane. Zero curvature preserves their parallelism indefinitely. These cases define the three canonical geometries in two dimensions: spherical, hyperbolic, and flat.

Curvature is a local property: it describes how space behaves in an infinitesimal neighbourhood. Compactness is a global property: it describes whether space is bounded and complete. The surface of a sphere is compact and positively curved. A flat plane is non-compact and uncurved. A cylinder is flat but compact in one direction. A torus has zero curvature but is compact. Curvature and compactness are independent notions.

A three-dimensional space can have its own curvature, defined purely through internal measurements of distance and angle, without requiring an external embedding. General relativity models the universe using such three-dimensional spatial geometries evolving in time.

The mathematical classification of homogeneous (uniform in all positions), isotropic (uniform in all directions) three-dimensional spaces yields three possibilities: positive curvature (a 3-sphere), negative curvature (a 3-hyperboloid), and zero curvature (Euclidean $\mathbb{R}^3$). Each corresponds to a constant value of spatial curvature and admits a well-defined metric. These are the geometric possibilities for the shape of the universe on large scales.

A cosmological model describes the evolution of space and time on the largest scales. In general relativity, such a model is not a visual rendering of stars and galaxies, but a mathematical solution to Einstein’s field equations (Einstein, 1915). It specifies the metric tensor—a geometric object encoding distances, angles, and causal relationships across spacetime. Given assumptions about symmetry and matter content, the metric determines how space stretches, curves, and evolves in time.

The standard model of cosmology is called the Lambda–Cold Dark Matter model ($\Lambda$CDM). It assumes that, at sufficiently large scales, the universe is homogeneous and isotropic, restricting the spatial geometry to one of the three cases described above. The spatial curvature in $\Lambda$CDM is not a free assumption. It is determined by the total energy density of the universe relative to a critical threshold. Density above, below, or at this value yields positive, negative, or zero curvature (Kolb \& Turner, 1990).

In this model, the Big Bang is not a point in space, but a boundary in time: a moment when the scale factor—the function describing the distance between any two comoving points—reaches zero. It represents the earliest definable state of the metric, beyond which classical general relativity ceases to apply. The Big Bang is not an explosion of matter into space. It is the dynamical expansion of space, or more accurately the distance between objects, governed by the evolving metric. All regions of the universe were arbitrarily close together in the past and have since expanded away from each other in a coordinated, metric-driven evolution.

This expansion is not centred at any specific location. It occurs everywhere simultaneously. Each observer sees distant galaxies receding, not because they are moving through space, but because the space between them is getting bigger. In a homogeneous universe, every region participates equally in the expansion, and the large-scale structure remains statistically uniform over time. The observable consequence is a redshift in the light from distant galaxies (Hubble, 1929)—a stretching of wavelengths that records the history of spatial expansion.

The strongest evidence for the geometry of the universe comes from the cosmic microwave background (CMB) (Penzias \& Wilson, 1965)—a relic radiation field that permeates all of space. The CMB originates from a time roughly 380,000 years after the Big Bang, when the universe cooled enough for protons and electrons to combine into neutral hydrogen atoms. This event, known as recombination, allowed photons to travel freely for the first time. The CMB is the redshifted remnant of that photon field, now observed at microwave wavelengths with a temperature of approximately 2.725 K.

The CMB is not perfectly uniform. It contains small anisotropies—tiny temperature fluctuations at the level of one part in 100,000. These fluctuations correspond to density variations in the early universe, which later seeded the formation of large-scale structures such as galaxies and clusters. The angular size of these fluctuations provides a direct measurement of spatial geometry. In particular, one can ask how large a primordial region appears on the sky today, given the time it took for light to reach us and the curvature of space through which it travelled.

Before recombination, the universe consisted of a hot, dense plasma of photons, electrons, and baryons (protons and neutrons). Photons scattered continuously off free electrons, coupling radiation tightly to matter. In this medium, density perturbations propagated as pressure waves driven by the competition between gravitational infall and photon pressure. These acoustic oscillations (Peebles \& Yu, 1970)—standing wave patterns in the plasma—left an imprint on the temperature distribution of the CMB when photon decoupling occurred.

The most important feature in the CMB spectrum is the first acoustic peak. This peak corresponds to the largest sound waves that had time to compress and rarefy a region of plasma before recombination. The physical size of such regions is determined by known physics—the speed of sound in the early universe and the duration before decoupling. However, the observed angular size depends on the spatial curvature of the universe. If space is positively curved, such a region appears larger than in flat geometry. If negatively curved, it appears smaller.

High-precision observations, particularly from the WMAP and Planck satellites, have measured the angular scale of the first acoustic peak with great accuracy. The result is consistent with a flat universe: the peak appears at an angle of approximately 1 degree, matching the prediction from a zero-curvature model. This agreement indicates that spatial curvature is consistent with zero to within a few parts per thousand.

Flat geometry means space neither curves back on itself nor terminates at a boundary (see \hyperref[ch:compacttwinparadox]{Chapter~14}). It continues indefinitely in every direction. If the universe is spatially flat everywhere, the simplest consistent model is one infinite in extent—no enclosing shell, no edge beyond which space ceases to exist. This conflicts with everyday intuition, which expects all things to be either bounded or looped. General relativity is indifferent to such expectations: it computes curvature from energy density and evolves the metric accordingly.

The observable universe does have a limit—the particle horizon, roughly 46 billion light-years in radius—but this is an observational boundary, not a physical wall. Beyond it, the model predicts space continues with the same statistical properties indefinitely.

That infinite space can arise from a finite beginning seems paradoxical, but time and space play different roles. The Big Bang marks a finite past moment—the origin of the metric and the expansion dynamics. If space was already infinite at that moment, it remained so as it expanded. The scale factor grew, but the global extent may always have been unbounded.

The universe's flatness is not a theoretical preference but an outcome of measurement. The cosmic microwave background, galaxy clustering, and large-scale structure all point to a flat metric, and the $\Lambda$CDM model achieves this without assuming it a priori.

The name \QSOpen{}$\Lambda$CDM\QSClose{} encodes the model's two dominant components: $\Lambda$ (Lambda) represents dark energy—the cosmological constant driving accelerated expansion (Riess et al., 1998)—and CDM stands for cold dark matter. Dark matter's role in structure formation is not incidental. During the radiation era, baryonic matter (protons, neutrons, electrons) remained tightly coupled to photons through electromagnetic interactions. Photon pressure prevented gravitational collapse. Dark matter, interacting only gravitationally, decoupled immediately after the Big Bang and began clustering in response to the density fluctuations seeded during inflation. By the time of recombination—when photons finally decoupled from matter—dark matter had already formed deep gravitational potential wells. Baryons then fell into these pre-existing halos, forming the first protogalaxies. Without this head start, the universe would not have had enough time to form the observed structures. The CMB fluctuations mentioned earlier are snapshots of these nascent dark matter concentrations. 

\begin{commentary}[ex nihilo]
The investigation of the origin of \QSOpen{}the world\QSClose{} is a central theme in all cultures—from spacetime to abiogenesis, between religion and science. Yet their answers are not mutually exclusive. Creation stories reveal nothing about the scientific account, and science could not care less about the intuitive question of why there is existence rather than nothing. It is interesting that the Big Bang model was initially ridiculed as too religious and later as atheistic ramblings, while in reality it's neither. It is simply the best set of equations that describes everything we observe to an unmatched degree of accuracy. The intuitive, not scientific, question of existence itself is in the domain of philosophy and theology, not science. The question of which model best explains the state of the universe, while being predictively fruitful, is not only in the \textit{domain} of science, but the \textit{raison d'être} of science.
\end{commentary}

\newpage
\thispagestyle{empty}

\begin{center}
\begin{tcolorbox}[
  enhanced,
  breakable,
  width=\textwidth,
  colframe=blue!60!black,
  colback=blue!5,
  colbacktitle=blue!60!black,
  coltitle=white,
  boxrule=0.5pt,
  arc=2mm,
  left=10pt,
  right=10pt,
  top=10pt,
  bottom=10pt,
  title=$\Lambda$CDM Cosmology Timeline,
  fonttitle=\bfseries\large,
  attach boxed title to top left={yshift=-2mm, xshift=5mm},
  boxed title style={arc=1mm, boxrule=0.5pt}
]

\begin{multicols}{2}
\raggedright
\footnotesize
\setlength{\parskip}{3pt}

{\bfseries 1. Planck Era ($<10^{-43}$ s)}\\
\textit{Forces:} Four forces unified. Quantum gravity dominates.\\
\textit{Matter:} No particles; quantum foam.\\
\textit{Scale:} Fluctuations at Planck length ($\sim10^{-35}$ m).

{\bfseries 2. Inflation ($10^{-36}$–$10^{-32}$ s)}\\
\textit{Forces:} Strong separates from electroweak.\\
\textit{Matter:} Vacuum energy drives exponential expansion.\\
\textit{Scale:} Universe expands by factor $\geq10^{26}$; quantum fluctuations seed galaxies.

{\bfseries 3. Reheating and Baryogenesis}\\
\textit{Time:} $10^{-32}$–$10^{-6}$ s\\
\textit{Forces:} Electroweak breaks into weak and electromagnetic.\\
\textit{Matter:} Quarks, leptons, gluons. CP violation creates 1:$10^9$ matter excess.\\
\textit{Temp:} $10^{15}$–$10^{12}$ K.

{\bfseries 4. Quark–Gluon Plasma}\\
\textit{Time:} $10^{-6}$–$10^{-5}$ s\\
\textit{Forces:} All four forces distinct.\\
\textit{Matter:} Quarks confine into protons and neutrons.\\
\textit{Temp:} $\sim10^{13}$ K.

{\bfseries 5. Nucleosynthesis (1 s–20 min)}\\
\textit{Matter:} Protons and neutrons fuse: 75\% H, 25\% He, trace Li.\\
\textit{Scale:} Photons coupled to matter; uniform plasma.

{\bfseries 6. Photon Era (20 min–380 kyr)}\\
\textit{Matter:} Ionised plasma. Neutrinos decouple at $\sim$1 s. CDM decouples immediately; non-interacting.\\
\textit{Temp:} $10^9$ K $\rightarrow$ 3000 K.\\
\textit{Scale:} Density fluctuations ($\sim10^{-5}$) grow. CDM begins clustering gravitationally.

\columnbreak

{\bfseries 7. Recombination and CMB}\\
\textit{Time:} 380,000 yr\\
\textit{Matter:} Electrons bind to nuclei $\rightarrow$ neutral atoms. CDM potential wells already formed.\\
\textit{Scale:} Photons decouple $\rightarrow$ CMB. Baryons begin falling into CDM halos.

{\bfseries 8. Dark Ages (0.4–0.5 Gyr)}\\
\textit{Matter:} Neutral hydrogen falls into CDM halos. CDM provides $\sim$85\% of gravitational scaffolding.\\
\textit{Scale:} Hierarchical merging; no stars yet. Smallest halos form first.

{\bfseries 9. First Stars and Reionization}\\
\textit{Time:} 0.5–1 Gyr\\
\textit{Matter:} Nuclear fusion creates first stars (Pop III); heavier elements forged.\\
\textit{Scale:} Ionising radiation clears fog; protogalaxies form.

{\bfseries 10. Galaxy Formation (1–5 Gyr)}\\
\textit{Scale:} CDM halos merge hierarchically; baryons cool and collapse at halo centres. Cosmic web emerges (100 Mpc filaments).\\
\textit{Activity:} Quasars peak ($\sim$3 Gyr). Galaxy rotation curves reveal CDM dominance.

{\bfseries 11. Present (13.8 Gyr)}\\
\textit{Forces:} Dark energy ($\Lambda$) dominates since $\sim$5–6 Gyr; drives accelerated expansion.\\
\textit{Composition:} 69\% dark energy, 26\% CDM, 5\% baryons.\\
\textit{Scale:} Expansion rate: $H_0 \approx 70$ km/s/Mpc. Galaxy separation $\sim$1 Mpc.\\
\textit{Effect:} $\Lambda$ prevents new large-scale structure formation; existing structures bound by CDM resist expansion.

{\bfseries 12. Future ($\Lambda$ domination)}\\
\textit{Near:} 10–100 Gyr: $\Lambda$-driven expansion pushes galaxies beyond event horizon. Only Local Group (bound by CDM) remains visible.\\
\textit{Mid:} $10^{12}$–$10^{14}$ yr: Star formation ceases; gas exhausted.\\
\textit{Far:} $10^{15}$–$10^{100}$ yr: Black holes evaporate. Universe asymptotes to cold, empty de Sitter space.

\end{multicols}

\vspace{4pt}
\hrule
\vspace{4pt}
{\footnotesize\textbf{Evolution Summary:} Unified forces $\rightarrow$ separate by $10^{-12}$ s. Quantum fields $\rightarrow$ quarks $\rightarrow$ hadrons $\rightarrow$ atoms $\rightarrow$ stars $\rightarrow$ galaxies. Planck scale $\rightarrow$ cosmic horizon. Radiation era $\rightarrow$ matter era (CDM-driven structure formation) $\rightarrow$ $\Lambda$ era (accelerated expansion).}

\end{tcolorbox}
\end{center}
